\IfFileExists{stacks-project.cls}{%
\documentclass{stacks-project}
}{%
\documentclass{amsart}
}

% For dealing with references we use the comment environment
\usepackage{verbatim}
\newenvironment{reference}{\comment}{\endcomment}
%\newenvironment{reference}{}{}
\newenvironment{slogan}{\comment}{\endcomment}
\newenvironment{history}{\comment}{\endcomment}

% For commutative diagrams we use Xy-pic
\usepackage[all]{xy}

% We use 2cell for 2-commutative diagrams.
\xyoption{2cell}
\UseAllTwocells

% We use multicol for the list of chapters between chapters
\usepackage{multicol}

% This is generally recommended for better output
\usepackage{lmodern}
\usepackage[T1]{fontenc}

% For cross-file-references
\usepackage{xr-hyper}

% Package for hypertext links:
\usepackage{hyperref}

% For any local file, say "hello.tex" you want to link to please
% use \externaldocument[hello-]{hello}
\externaldocument[introduction-]{introduction}
\externaldocument[conventions-]{conventions}
\externaldocument[sets-]{sets}
\externaldocument[categories-]{categories}
\externaldocument[topology-]{topology}
\externaldocument[sheaves-]{sheaves}
\externaldocument[sites-]{sites}
\externaldocument[stacks-]{stacks}
\externaldocument[fields-]{fields}
\externaldocument[algebra-]{algebra}
\externaldocument[brauer-]{brauer}
\externaldocument[homology-]{homology}
\externaldocument[derived-]{derived}
\externaldocument[simplicial-]{simplicial}
\externaldocument[more-algebra-]{more-algebra}
\externaldocument[smoothing-]{smoothing}
\externaldocument[modules-]{modules}
\externaldocument[sites-modules-]{sites-modules}
\externaldocument[injectives-]{injectives}
\externaldocument[cohomology-]{cohomology}
\externaldocument[sites-cohomology-]{sites-cohomology}
\externaldocument[dga-]{dga}
\externaldocument[dpa-]{dpa}
\externaldocument[sdga-]{sdga}
\externaldocument[hypercovering-]{hypercovering}
\externaldocument[schemes-]{schemes}
\externaldocument[constructions-]{constructions}
\externaldocument[properties-]{properties}
\externaldocument[morphisms-]{morphisms}
\externaldocument[coherent-]{coherent}
\externaldocument[divisors-]{divisors}
\externaldocument[limits-]{limits}
\externaldocument[varieties-]{varieties}
\externaldocument[topologies-]{topologies}
\externaldocument[descent-]{descent}
\externaldocument[perfect-]{perfect}
\externaldocument[more-morphisms-]{more-morphisms}
\externaldocument[flat-]{flat}
\externaldocument[groupoids-]{groupoids}
\externaldocument[more-groupoids-]{more-groupoids}
\externaldocument[etale-]{etale}
\externaldocument[chow-]{chow}
\externaldocument[intersection-]{intersection}
\externaldocument[pic-]{pic}
\externaldocument[weil-]{weil}
\externaldocument[adequate-]{adequate}
\externaldocument[dualizing-]{dualizing}
\externaldocument[duality-]{duality}
\externaldocument[discriminant-]{discriminant}
\externaldocument[derham-]{derham}
\externaldocument[local-cohomology-]{local-cohomology}
\externaldocument[algebraization-]{algebraization}
\externaldocument[curves-]{curves}
\externaldocument[resolve-]{resolve}
\externaldocument[models-]{models}
\externaldocument[functors-]{functors}
\externaldocument[equiv-]{equiv}
\externaldocument[pione-]{pione}
\externaldocument[etale-cohomology-]{etale-cohomology}
\externaldocument[proetale-]{proetale}
\externaldocument[relative-cycles-]{relative-cycles}
\externaldocument[more-etale-]{more-etale}
\externaldocument[trace-]{trace}
\externaldocument[crystalline-]{crystalline}
\externaldocument[spaces-]{spaces}
\externaldocument[spaces-properties-]{spaces-properties}
\externaldocument[spaces-morphisms-]{spaces-morphisms}
\externaldocument[decent-spaces-]{decent-spaces}
\externaldocument[spaces-cohomology-]{spaces-cohomology}
\externaldocument[spaces-limits-]{spaces-limits}
\externaldocument[spaces-divisors-]{spaces-divisors}
\externaldocument[spaces-over-fields-]{spaces-over-fields}
\externaldocument[spaces-topologies-]{spaces-topologies}
\externaldocument[spaces-descent-]{spaces-descent}
\externaldocument[spaces-perfect-]{spaces-perfect}
\externaldocument[spaces-more-morphisms-]{spaces-more-morphisms}
\externaldocument[spaces-flat-]{spaces-flat}
\externaldocument[spaces-groupoids-]{spaces-groupoids}
\externaldocument[spaces-more-groupoids-]{spaces-more-groupoids}
\externaldocument[bootstrap-]{bootstrap}
\externaldocument[spaces-pushouts-]{spaces-pushouts}
\externaldocument[spaces-chow-]{spaces-chow}
\externaldocument[groupoids-quotients-]{groupoids-quotients}
\externaldocument[spaces-more-cohomology-]{spaces-more-cohomology}
\externaldocument[spaces-simplicial-]{spaces-simplicial}
\externaldocument[spaces-duality-]{spaces-duality}
\externaldocument[formal-spaces-]{formal-spaces}
\externaldocument[restricted-]{restricted}
\externaldocument[spaces-resolve-]{spaces-resolve}
\externaldocument[formal-defos-]{formal-defos}
\externaldocument[defos-]{defos}
\externaldocument[cotangent-]{cotangent}
\externaldocument[examples-defos-]{examples-defos}
\externaldocument[algebraic-]{algebraic}
\externaldocument[examples-stacks-]{examples-stacks}
\externaldocument[stacks-sheaves-]{stacks-sheaves}
\externaldocument[criteria-]{criteria}
\externaldocument[artin-]{artin}
\externaldocument[quot-]{quot}
\externaldocument[stacks-properties-]{stacks-properties}
\externaldocument[stacks-morphisms-]{stacks-morphisms}
\externaldocument[stacks-limits-]{stacks-limits}
\externaldocument[stacks-cohomology-]{stacks-cohomology}
\externaldocument[stacks-perfect-]{stacks-perfect}
\externaldocument[stacks-introduction-]{stacks-introduction}
\externaldocument[stacks-more-morphisms-]{stacks-more-morphisms}
\externaldocument[stacks-geometry-]{stacks-geometry}
\externaldocument[moduli-]{moduli}
\externaldocument[moduli-curves-]{moduli-curves}
\externaldocument[examples-]{examples}
\externaldocument[exercises-]{exercises}
\externaldocument[guide-]{guide}
\externaldocument[desirables-]{desirables}
\externaldocument[coding-]{coding}
\externaldocument[obsolete-]{obsolete}
\externaldocument[fdl-]{fdl}
\externaldocument[index-]{index}

% Theorem environments.
%
\theoremstyle{plain}
\newtheorem{theorem}[subsection]{Theorem}
\newtheorem{proposition}[subsection]{Proposition}
\newtheorem{lemma}[subsection]{Lemma}

\theoremstyle{definition}
\newtheorem{definition}[subsection]{Definition}
\newtheorem{example}[subsection]{Example}
\newtheorem{exercise}[subsection]{Exercise}
\newtheorem{situation}[subsection]{Situation}

\theoremstyle{remark}
\newtheorem{remark}[subsection]{Remark}
\newtheorem{remarks}[subsection]{Remarks}

\numberwithin{equation}{subsection}

% Macros
%
\def\lim{\mathop{\mathrm{lim}}\nolimits}
\def\colim{\mathop{\mathrm{colim}}\nolimits}
\def\Spec{\mathop{\mathrm{Spec}}}
\def\Hom{\mathop{\mathrm{Hom}}\nolimits}
\def\Ext{\mathop{\mathrm{Ext}}\nolimits}
\def\SheafHom{\mathop{\mathcal{H}\!\mathit{om}}\nolimits}
\def\SheafExt{\mathop{\mathcal{E}\!\mathit{xt}}\nolimits}
\def\Sch{\mathit{Sch}}
\def\Mor{\mathop{\mathrm{Mor}}\nolimits}
\def\Ob{\mathop{\mathrm{Ob}}\nolimits}
\def\Sh{\mathop{\mathit{Sh}}\nolimits}
\def\NL{\mathop{N\!L}\nolimits}
\def\CH{\mathop{\mathrm{CH}}\nolimits}
\def\proetale{{pro\text{-}\acute{e}tale}}
\def\etale{{\acute{e}tale}}
\def\QCoh{\mathit{QCoh}}
\def\Ker{\mathop{\mathrm{Ker}}}
\def\Im{\mathop{\mathrm{Im}}}
\def\Coker{\mathop{\mathrm{Coker}}}
\def\Coim{\mathop{\mathrm{Coim}}}

% Boxtimes
%
\DeclareMathSymbol{\boxtimes}{\mathbin}{AMSa}{"02}

%
% Macros for moduli stacks/spaces
%
\def\QCohstack{\mathcal{QC}\!\mathit{oh}}
\def\Cohstack{\mathcal{C}\!\mathit{oh}}
\def\Spacesstack{\mathcal{S}\!\mathit{paces}}
\def\Quotfunctor{\mathrm{Quot}}
\def\Hilbfunctor{\mathrm{Hilb}}
\def\Curvesstack{\mathcal{C}\!\mathit{urves}}
\def\Polarizedstack{\mathcal{P}\!\mathit{olarized}}
\def\Complexesstack{\mathcal{C}\!\mathit{omplexes}}
% \Pic is the operator that assigns to X its picard group, usage \Pic(X)
% \Picardstack_{X/B} denotes the Picard stack of X over B
% \Picardfunctor_{X/B} denotes the Picard functor of X over B
\def\Pic{\mathop{\mathrm{Pic}}\nolimits}
\def\Picardstack{\mathcal{P}\!\mathit{ic}}
\def\Picardfunctor{\mathrm{Pic}}
\def\Deformationcategory{\mathcal{D}\!\mathit{ef}}
\usepackage{cleveref}
\crefname{originalequation}{original equation}{original equations}
\Crefname{originalequation}{Original equation}{Original equations}
\makeatletter
\newlabel{equation-formal-sum}{{0F6J}{source}{}{sourcepointer-1}{}}
\newlabel{equation-formal-sum@cref}{{[originalequation][0][]0F6J}{[1][0][]source}{}{}{}}
\makeatother
\crefname{originalremark}{original remark}{original remarks}
\Crefname{originalremark}{Original remark}{Original remarks}
\makeatletter
\newlabel{etale-cohomology-remark-stalk-pullback}{{04JN}{source}{}{sourcepointer-2}{}}
\newlabel{etale-cohomology-remark-stalk-pullback@cref}{{[originalremark][0][]04JN}{[1][0][]source}{}{}{}}
\makeatother
\crefname{originaltheorem}{original theorem}{original theorems}
\Crefname{originaltheorem}{Original theorem}{Original theorems}
\makeatletter
\newlabel{etale-cohomology-theorem-exactness-stalks}{{03PU}{source}{}{sourcepointer-3}{}}
\newlabel{etale-cohomology-theorem-exactness-stalks@cref}{{[originaltheorem][0][]03PU}{[1][0][]source}{}{}{}}
\makeatother
\crefname{originalreference}{original reference}{original references}
\Crefname{originalreference}{Original reference}{Original references}
\makeatletter
\newlabel{item-sub}{{0F6L}{source}{}{sourcepointer-4}{}}
\newlabel{item-sub@cref}{{[originalreference][0][]0F6L}{[1][0][]source}{}{}{}}
\makeatother
\crefname{originalreference}{original reference}{original references}
\Crefname{originalreference}{Original reference}{Original references}
\makeatletter
\newlabel{item-sum}{{0F6K}{source}{}{sourcepointer-5}{}}
\newlabel{item-sum@cref}{{[originalreference][0][]0F6K}{[1][0][]source}{}{}{}}
\makeatother
\crefname{originallemma}{original lemma}{original lemmas}
\Crefname{originallemma}{Original lemma}{Original lemmas}
\makeatletter
\newlabel{lemma-base-change-f-shriek-separated}{{0F55}{source}{}{sourcepointer-6}{}}
\newlabel{lemma-base-change-f-shriek-separated@cref}{{[originallemma][0][]0F55}{[1][0][]source}{}{}{}}
\makeatother
\crefname{originallemma}{original lemma}{original lemmas}
\Crefname{originallemma}{Original lemma}{Original lemmas}
\makeatletter
\newlabel{lemma-f-shriek-composition}{{0F50}{source}{}{sourcepointer-7}{}}
\newlabel{lemma-f-shriek-composition@cref}{{[originallemma][0][]0F50}{[1][0][]source}{}{}{}}
\makeatother
\crefname{originallemma}{original lemma}{original lemmas}
\Crefname{originallemma}{Original lemma}{Original lemmas}
\makeatletter
\newlabel{lemma-finite-support-f-shriek-separated}{{0F6N}{source}{}{sourcepointer-8}{}}
\newlabel{lemma-finite-support-f-shriek-separated@cref}{{[originallemma][0][]0F6N}{[1][0][]source}{}{}{}}
\makeatother
\crefname{originallemma}{original lemma}{original lemmas}
\Crefname{originallemma}{Original lemma}{Original lemmas}
\makeatletter
\newlabel{lemma-lqf-base-change-f-shriek}{{0F5J}{source}{}{sourcepointer-9}{}}
\newlabel{lemma-lqf-base-change-f-shriek@cref}{{[originallemma][0][]0F5J}{[1][0][]source}{}{}{}}
\makeatother
\crefname{originallemma}{original lemma}{original lemmas}
\Crefname{originallemma}{Original lemma}{Original lemmas}
\makeatletter
\newlabel{lemma-lqf-f-shriek-stalk}{{0F5F}{source}{}{sourcepointer-10}{}}
\newlabel{lemma-lqf-f-shriek-stalk@cref}{{[originallemma][0][]0F5F}{[1][0][]source}{}{}{}}
\makeatother
\crefname{originallemma}{original lemma}{original lemmas}
\Crefname{originallemma}{Original lemma}{Original lemmas}
\makeatletter
\newlabel{lemma-stalk-f-shriek-separated}{{0F76}{source}{}{sourcepointer-11}{}}
\newlabel{lemma-stalk-f-shriek-separated@cref}{{[originallemma][0][]0F76}{[1][0][]source}{}{}{}}
\makeatother
\crefname{originallemma}{original lemma}{original lemmas}
\Crefname{originallemma}{Original lemma}{Original lemmas}
\makeatletter
\newlabel{morphisms-lemma-finite-permanence}{{035D}{source}{}{sourcepointer-12}{}}
\newlabel{morphisms-lemma-finite-permanence@cref}{{[originallemma][0][]035D}{[1][0][]source}{}{}{}}
\makeatother
\crefname{originallemma}{original lemma}{original lemmas}
\Crefname{originallemma}{Original lemma}{Original lemmas}
\makeatletter
\newlabel{morphisms-lemma-scheme-theoretic-image-is-proper}{{0AH6}{source}{}{sourcepointer-13}{}}
\newlabel{morphisms-lemma-scheme-theoretic-image-is-proper@cref}{{[originallemma][0][]0AH6}{[1][0][]source}{}{}{}}
\makeatother
\crefname{originalremark}{original remark}{original remarks}
\Crefname{originalremark}{Original remark}{Original remarks}
\makeatletter
\newlabel{remark-construct-map-presheaves-downstairs}{{0F78}{source}{}{sourcepointer-14}{}}
\newlabel{remark-construct-map-presheaves-downstairs@cref}{{[originalremark][0][]0F78}{[1][0][]source}{}{}{}}
\makeatother
\crefname{originalremark}{original remark}{original remarks}
\Crefname{originalremark}{Original remark}{Original remarks}
\makeatletter
\newlabel{remark-covariance-f-shriek-separated}{{0F53}{source}{}{sourcepointer-15}{}}
\newlabel{remark-covariance-f-shriek-separated@cref}{{[originalremark][0][]0F53}{[1][0][]source}{}{}{}}
\makeatother
\crefname{originalremark}{original remark}{original remarks}
\Crefname{originalremark}{Original remark}{Original remarks}
\makeatletter
\newlabel{remark-covariance-lqf-f-shriek}{{0F6S}{source}{}{sourcepointer-16}{}}
\newlabel{remark-covariance-lqf-f-shriek@cref}{{[originalremark][0][]0F6S}{[1][0][]source}{}{}{}}
\makeatother
\crefname{originallemma}{original lemma}{original lemmas}
\Crefname{originallemma}{Original lemma}{Original lemmas}
\makeatletter
\newlabel{varieties-lemma-algebraic-scheme-dim-0}{{06LH}{source}{}{sourcepointer-17}{}}
\newlabel{varieties-lemma-algebraic-scheme-dim-0@cref}{{[originallemma][0][]06LH}{[1][0][]source}{}{}{}}
\makeatother
\title{Composition isomorphism g\_! composed with f'\_! equals f\_! for locally quasi-finite f' and separated locally finite type g}
\author{The Stacks Project Authors}
\date{stacks-project:more-etale.tex:local-snapshot-20261012; GNU FDL 1.2 or later}
\begin{document}
\maketitle
\setcounter{subsection}{1}
\noindent\textit{Editorial note.} All delivered mathematical text is the author's own, copied from the pinned Stacks Project chapter file more-etale.tex; no statement or proof step was rewritten, reordered or removed. The counted span is the author environment and its own proof environment. References inside the span are handed to the assembler as navigation stubs that carry the canonical tag of the cited tag only; whether each of them is load-bearing is recorded on the stub itself.\par\medskip
\noindent\textit{Editorial note.} Disclosure: part (c) of the statement (compatibility with Lemma lemma-f-shriek-composition via Lemma lemma-finite-support-f-shriek-separated when f' is separated) is proved by the author only as 'This follows from comparing the definitions for both maps; we omit the details.' (proof lines 1791-1792). The canonical isomorphism itself and parts (a) and (b) are fully proved inside the counted span; part (c) is an ancillary compatibility not used in the isomorphism argument.\par\medskip
\noindent\textit{Allowed standard prerequisites (editorial).}\begin{itemize}
\item Presheaf-map construction conventions for lower shriek (Remark remark-construct-map-presheaves-downstairs, canonical tag 0F78). cited at proof lines 1643-1651. Fixes the conventions under which the map c determines the canonical map c\textasciicircum{}\# on sheafifications.
\item Exactness of stalk functors (Etale Cohomology, Theorem etale-cohomology-theorem-exactness-stalks, canonical tag 03PU). cited at proof lines 1800-1801. Justifies proving the isomorphism c' on stalks.
\end{itemize}\medskip
\ifdefined\hypertarget\hypertarget{sourcepointer-1}{}\fi
\noindent\textit{Source reference.} Original equation 0F6J: sources/stacks-chow/tags.txt, source lines 18974--18974. Cited at proof line 1655: the formal-sum presentation convention that the map c acts on.\par\medskip
\ifdefined\hypertarget\hypertarget{sourcepointer-2}{}\fi
\noindent\textit{Source reference.} Original remark 04JN: sources/stacks-chow/tags.txt, source lines 5944--5944. Cited at proof line 1799: pulling back by the geometric point is taking stalks.\par\medskip
\ifdefined\hypertarget\hypertarget{sourcepointer-3}{}\fi
\noindent\textit{Source reference.} Original theorem 03PU: sources/stacks-chow/tags.txt, source lines 4871--4871. Cited at proof line 1801: exactness of stalk functors lets the isomorphism be checked on stalks.\par\medskip
\ifdefined\hypertarget\hypertarget{sourcepointer-4}{}\fi
\noindent\textit{Source reference.} Original reference 0F6L: sources/stacks-chow/tags.txt, source lines 18976--18976. Cited at proof line 1663: the map c is compatible with relation (sub).\par\medskip
\ifdefined\hypertarget\hypertarget{sourcepointer-5}{}\fi
\noindent\textit{Source reference.} Original reference 0F6K: sources/stacks-chow/tags.txt, source lines 18975--18975. Cited at proof line 1663: the map c is compatible with relation (sum).\par\medskip
\ifdefined\hypertarget\hypertarget{sourcepointer-6}{}\fi
\noindent\textit{Source reference.} Original lemma 0F55: sources/stacks-chow/tags.txt, source lines 18925--18925. Cited at statement line 1630: base change isomorphisms of Lemma (b), separated case.\par\medskip
\ifdefined\hypertarget\hypertarget{sourcepointer-7}{}\fi
\noindent\textit{Source reference.} Original lemma 0F50: sources/stacks-chow/tags.txt, source lines 18920--18920. Cited at statement line 1632: the isomorphism that part (c) compares with.\par\medskip
\ifdefined\hypertarget\hypertarget{sourcepointer-8}{}\fi
\noindent\textit{Source reference.} Original lemma 0F6N: sources/stacks-chow/tags.txt, source lines 18978--18978. Cited at statement line 1633: the identifications used by part (c) when f' is separated.\par\medskip
\ifdefined\hypertarget\hypertarget{sourcepointer-9}{}\fi
\noindent\textit{Source reference.} Original lemma 0F5J: sources/stacks-chow/tags.txt, source lines 18939--18939. Cited at statement line 1629: base change isomorphisms of Lemma (b), locally quasi-finite case.\par\medskip
\ifdefined\hypertarget\hypertarget{sourcepointer-10}{}\fi
\noindent\textit{Source reference.} Original lemma 0F5F: sources/stacks-chow/tags.txt, source lines 18935--18935. Cited at proof lines 1812 and 1821: stalk of f'\_!F and the skyscraper sheaf at the image point.\par\medskip
\ifdefined\hypertarget\hypertarget{sourcepointer-11}{}\fi
\noindent\textit{Source reference.} Original lemma 0F76: sources/stacks-chow/tags.txt, source lines 18996--18996. Cited at proof line 1813: the H\textasciicircum{}0\_c identification used in the isomorphism check.\par\medskip
\ifdefined\hypertarget\hypertarget{sourcepointer-12}{}\fi
\noindent\textit{Source reference.} Original lemma 035D: sources/stacks-chow/tags.txt, source lines 4124--4124. Cited at proof line 1661: Z\_i is finite over V' because g is separated.\par\medskip
\ifdefined\hypertarget\hypertarget{sourcepointer-13}{}\fi
\noindent\textit{Source reference.} Original lemma 0AH6: sources/stacks-chow/tags.txt, source lines 13219--13219. Cited at proof lines 1672-1673: the image of the disjoint union of the Z\_i is proper over V.\par\medskip
\ifdefined\hypertarget\hypertarget{sourcepointer-14}{}\fi
\noindent\textit{Source reference.} Original remark 0F78: sources/stacks-chow/tags.txt, source lines 18998--18998. Cited at proof lines 1643-1651: the conventions under which c determines the canonical c\textasciicircum{}\#.\par\medskip
\ifdefined\hypertarget\hypertarget{sourcepointer-15}{}\fi
\noindent\textit{Source reference.} Original remark 0F53: sources/stacks-chow/tags.txt, source lines 18923--18923. Cited at statement line 1627: covariance compatibility (a), separated case.\par\medskip
\ifdefined\hypertarget\hypertarget{sourcepointer-16}{}\fi
\noindent\textit{Source reference.} Original remark 0F6S: sources/stacks-chow/tags.txt, source lines 18982--18982. Cited at statement line 1628: covariance compatibility (a), locally quasi-finite case.\par\medskip
\ifdefined\hypertarget\hypertarget{sourcepointer-17}{}\fi
\noindent\textit{Source reference.} Original lemma 06LH: sources/stacks-chow/tags.txt, source lines 8466--8466. Cited at proof line 1816: X is a disjoint union of spectra of Artinian local rings with residue field k.\par\medskip
\begin{lemma}
\label{lemma-lqf-separated-shriek-composition}
Let $f' : X \to Y'$ and $g : Y' \to Y$ be composable morphisms of schemes
with $f'$ and $f = g \circ f'$ locally quasi-finite and $g$ separated and
locally of finite type. Then there is a canonical isomorphism of functors
$g_! \circ f'_! = f_!$. This isomorphism is compatible with
\begin{enumerate}
\item[(a)] covariance with respect to open embeddings as in
Remarks \ref{remark-covariance-f-shriek-separated} and
\ref{remark-covariance-lqf-f-shriek},
\item[(b)] the base change isomorphisms of
Lemmas \ref{lemma-lqf-base-change-f-shriek}
and \ref{lemma-base-change-f-shriek-separated}, and
\item[(c)] equal to the isomorphism of Lemma \ref{lemma-f-shriek-composition}
via the identifications of Lemma \ref{lemma-finite-support-f-shriek-separated}
in case $f'$ is separated.
\end{enumerate}
\end{lemma}

\begin{proof}
Let $\mathcal{F}$ be an abelian sheaf on $X_\etale$. With conventions as in
Remark \ref{remark-construct-map-presheaves-downstairs} we will explicitly
construct a map
$$
c : f_{p!}\mathcal{F} \longrightarrow g_*f'_{p!}\mathcal{F}
$$
of abelian presheaves on $Y_\etale$. By the discussion in
Remark \ref{remark-construct-map-presheaves-downstairs}
this will determine a canonical map
$c^\# : f_!\mathcal{F} \to g_*f'_!\mathcal{F}$.
We will show that $c^\#$ has image contained in the subsheaf
$g_!f'_!\mathcal{F}$, thereby obtaining a map
$c' : f_!\mathcal{F} \to g_!f'_!\mathcal{F}$. Next, we will prove
(a), (b), and (c) that. Finally, part (b)
will allow us to show that $c'$ is an isomorphism.

\medskip\noindent
Construction of the map $c$. Let $V \in Y_\etale$ and
let $s = \sum (Z_i, s_i)$ be a sum as in (\ref{equation-formal-sum})
defining an element of $f_{p!}\mathcal{F}(V)$.
Recall that $Z_i \subset X_V = X \times_Y V$
is a locally closed subscheme finite over $V$.
Setting $V' = Y' \times_Y V$ we get $X_{V'} = X \times_{Y'} V' = X_V$.
Hence $Z_i \subset X_{V'}$ is locally closed and
$Z_i$ is finite over $V'$ because $g$ is separated
(Morphisms, Lemma \ref{morphisms-lemma-finite-permanence}).
Hence we may set $c(s) = \sum (Z_i, s_i)$ but now viewed
as an element of $f'_{p!}\mathcal{F}(V') = (g_*f'_{p!}\mathcal{F})(V)$.
The construction is clearly compatible with relations
(\ref{item-sum}) and (\ref{item-sub})
and compatible with restriction mappings and hence we obtain the map $c$.

\medskip\noindent
Observe that in the discussion above our section $c(s) = \sum (Z_i, s_i)$ of
$f'_!\mathcal{F}$ over $V'$ restricts to zero on
$V' \setminus \Im(\coprod Z_i \to V')$. Since $\Im(\coprod Z_i \to V')$
is proper over $V$ (for example by Morphisms, Lemma
\ref{morphisms-lemma-scheme-theoretic-image-is-proper})
we conclude that $c(s)$ defines a section of
$g_!f'_!\mathcal{F} \subset g_*f'_!\mathcal{F}$ over $V$.
Since every local section of $f_!\mathcal{F}$ locally comes from a
local section of $f_{p!}\mathcal{F}$ we conclude that the image
of $c^\#$ is contained in $g_!f'_!\mathcal{F}$.
Thus we obtain an induced map $c' : f_!\mathcal{F} \to g_!f'_!\mathcal{F}$
factoring $c^\#$ as predicted in the first paragraph of the proof.

\medskip\noindent
Proof of (a). Let $Y'_1 \subset Y'$ be an open subscheme
and set $X_1 = (f')^{-1}(W')$. We obtain a diagram
$$
\xymatrix{
X_1 \ar[d]_{f'_1} \ar[r]_a \ar@/_2em/[dd]_{f_1} &
X \ar[d]^{f'} \ar@/^2em/[dd]^f \\
Y'_1 \ar[d]_{g_1} \ar[r]_{b'} &
Y' \ar[d]^g \\
Y \ar@{=}[r] &
Y
}
$$
where the horizontal arrows are open immersions. Then our claim is that
the diagram
$$
\xymatrix{
f_{1, !}\mathcal{F}|_{X_1} \ar[r]_{c'_1} \ar[dd] &
g_{1, !}f'_{1, !}\mathcal{F}|_{X_1} \ar@{=}[d] \\
& g_{1, !}(f'_!\mathcal{F})|_{Y'_1} \ar[d] \\
f_!\mathcal{F} \ar[r]^{c'} &
g_!f'_!\mathcal{F} \ar[r] & g_*f'_!\mathcal{F}
}
$$
commutes where the left vertical arrow is
Remark \ref{remark-covariance-lqf-f-shriek} and
the right vertical arrow is Remark \ref{remark-covariance-f-shriek-separated}.
The equality sign in the diagram comes about because $f'_1$
is the restriction of $f'$ to $Y'_1$ and our construction
of $f'_!$ is local on the base.
Finally, to prove the commutativity we choose an object $V$ of
$Y_\etale$ and a formal sum $s_1 = \sum (Z_{1, i}, s_{1, i})$ as in
(\ref{equation-formal-sum}) defining an element of
$f_{1, p!}\mathcal{F}|_{X_1}(V)$. Recall this means
$Z_{1, i} \subset X_1 \times_Y V$ is locally closed finite over $V$
and $s_{1, i} \in H_{Z_{1, i}}(\mathcal{F})$.
Then we chase this section
across the maps involved, but we only need to show we
end up with the same element of
$g_*f'_!\mathcal{F}(V) = f'_!\mathcal{F}(Y' \times_Y V)$.
Going around both sides of the diagram the reader immediately
sees we end up with the element $\sum (Z_{1, i}, s_{1, i})$
where now $Z_{1, i}$ is viewed as a locally closed subscheme
of $X \times_{Y'} (Y' \times_Y V) = X \times_Y V$ finite over
$Y' \times_Y V$.

\medskip\noindent
Proof of (b). Let $b : Y_1 \to Y$ be a morphism of schemes. Let us form the
commutative diagram
$$
\xymatrix{
X_1 \ar[d]_{f'_1} \ar[r]_a \ar@/_2em/[dd]_{f_1} &
X \ar[d]^{f'} \ar@/^2em/[dd]^f \\
Y'_1 \ar[d]_{g_1} \ar[r]_{b'} &
Y' \ar[d]^g \\
Y_1 \ar[r]^b &
Y
}
$$
with cartesian squares. We claim that our construction is compatible
with the base change maps of Lemmas \ref{lemma-lqf-base-change-f-shriek}
and \ref{lemma-base-change-f-shriek-separated}, i.e.,
that the top rectangle of the diagram
$$
\xymatrix{
b^{-1}f_!\mathcal{F} \ar[rr] \ar[d]_{b^{-1}c'} & &
f_{1, !}a^{-1}\mathcal{F} \ar[d]^{c_1'} \\
b^{-1}g_!f'_!\mathcal{F} \ar[r] \ar[d] &
g_{1, !}(b')^{-1}f'_!\mathcal{F} \ar[r] \ar[d] &
g_{1, !}f'_{1, !}a^{-1}\mathcal{F} \ar[d] \\
b^{-1}g_*f'_!\mathcal{F} \ar[r] &
g_{1, *}(b')^{-1}f'_!\mathcal{F} \ar[r] &
g_{1, *}f'_{1, !}a^{-1}\mathcal{F}
}
$$
commutes. The verification of this is completely routine and we
urge the reader to skip it. Since the arrows going from the middle
row down to the bottom row are injective, it suffices to show that
the outer diagram commutes.
To show this it suffices to take a local section of
$b^{-1}f_!\mathcal{F}$ and show we end up with the same local
section of $g_{1, *}f'_{1, !}a^{-1}\mathcal{F}$
going around either way. However, in fact it suffices to check
this for local sections which are of the pullback by $b$ of
a section $s = \sum (Z_i, s_i)$ of $f_{p!}\mathcal{F}(V)$
as above (since such pullbacks generate the abelian sheaf
$b^{-1}f_!\mathcal{F}$). Denote $V_1$, $V'_1$, and $Z_{1, i}$
the base change of $V$, $V' = Y' \times_Y V$, $Z_i$ by $Y_1 \to Y$.
Recall that $Z_i$ is a locally closed subscheme of $X_V = X_{V'}$
and hence $Z_{1, i}$ is a locally closed subscheme
of $(X_1)_{V_1} = (X_1)_{V'_1}$. Then $b^{-1}c'$ sends the pullback
of $s$ to the pullback of the local section $c(s) \sum (Z_i, s_i)$ viewed
as an element of $f'_{p!}\mathcal{F}(V') = (g_*f'_{p!}\mathcal{F})(V)$.
The composition of the bottom two base change maps
simply maps this to $\sum (Z_{i, 1}, s_{1, i})$ viewed as an
element of $f'_{1, p!}a^{-1}\mathcal{F}(V'_1) =
g_{1, *}f'_{1, p!}a^{-1}\mathcal{F}(V_1)$.
On the other hand, the base change map at the top of the diagram
sends the pullback of $s$ to $\sum (Z_{1, i}, s_{1, i})$ viewed
as an element of $f_{1, !}a^{-1}\mathcal{F}(V_1)$.
Then finally $c'_1$ by its very construction does indeed
map this to $\sum (Z_{i, 1}, s_{1, i})$ viewed as an
element of $f'_{1, p!}a^{-1}\mathcal{F}(V'_1) =
g_{1, *}f'_{1, p!}a^{-1}\mathcal{F}(V_1)$ and the commutativity
has been verified.

\medskip\noindent
Proof of (c). This follows from comparing the definitions
for both maps; we omit the details.

\medskip\noindent
To finish the proof it suffices to show that the pullback of
$c'$ via any geometric point $\overline{y} : \Spec(k) \to Y$
is an isomorphism. Namely, pulling back by $\overline{y}$
is the same thing as taking stalks and $\overline{y}$
(\'Etale Cohomology, Remark \ref{etale-cohomology-remark-stalk-pullback})
and hence we can invoke
\'Etale Cohomology, Theorem \ref{etale-cohomology-theorem-exactness-stalks}.
By the compatibility (b) just shown, we 
conclude that we may assume $Y$ is the spectrum of $k$
and we have to show that $c'$ is an isomorphism.
To do this it suffices to show that the induced map
$$
\bigoplus\nolimits_{x \in X} \mathcal{F}_x = H^0(Y, f_!\mathcal{F})
\longrightarrow
H^0(Y, g_!f'_!\mathcal{F}) = H^0_c(Y', f'_!\mathcal{F})
$$
is an isomorphism. The equalities hold by
Lemmas \ref{lemma-lqf-f-shriek-stalk} and
\ref{lemma-stalk-f-shriek-separated}.
Recall that $X$ is a disjoint union of
spectra of Artinian local rings with residue field $k$, see
Varieties, Lemma \ref{varieties-lemma-algebraic-scheme-dim-0}.
Since the left and right hand side commute with direct
sums (details omitted) we may assume that $\mathcal{F}$ is a skyscraper
sheaf $x_*A$ supported at some $x \in X$.
Then $f'_!\mathcal{F}$ is the skyscraper sheaf at the
image $y'$ of $x$ in $Y$ by Lemma \ref{lemma-lqf-f-shriek-stalk}.
In this case it is obvious that our construction
produces the identity map $A \to H^0_c(Y', y'_*A) = A$
as desired.
\end{proof}

\end{document}
